In the same session, we varied retrieval depth, reranking depth and token limits. At 192 tokens, each RRF candidate setting was tested with no reranking and with TinyBERT reranking up to 8, 24 or 50 documents. Each suffix uses the full shared retrieval time plus its own timed reranking. At $K=20$, the candidate union has at most 40 documents. Raising per-retriever $K$ across 20, 64, and 100 increased candidate recall on all three domains (Table~\ref{tab:depth}), yet RRF NDCG@10 shifted at most 0.0031 from $K=64$ to $K=100$ (intervals crossing zero). At $K=64$, extending TinyBERT from 24 to 50 lowered NDCG@10 in all three domains by point estimate and added 23.6--31.7 ms. The SciFact interval includes zero. The NFCorpus and FiQA intervals are below zero.

At $K=64$ and reranking depth 24, we varied the encoder and reranker limits independently from 192 to 384 tokens. No query exceeded 192 tokens, whereas 94.9\%, 90.8\%, and 48.5\% of top-24 pairs did on SciFact, NFCorpus, and FiQA (89.0\%, 90.7\%, and 28.9\% of documents). The 192-token dense index is V5; 384-token encoding rebuilt each corpus offline (9.5, 6.9, and 93.7 s). Encoder-only expansion left NDCG@10 flat (intervals crossing zero). Reranker-only expansion raised NDCG@10 on SciFact by +0.0478 [0.0260, 0.0702] (+6.12 ms) and NFCorpus by +0.0083 [0.0022, 0.0148] (+5.27 ms), with uncertainty on FiQA (+0.0057 [-0.0027, 0.0141], +6.24 ms). At the middle budget, longer reranking raised utility only on SciFact (0.5494 to 0.5594), lowering NFCorpus (0.2630 to 0.2460) and FiQA (0.1300 to 0.0660). Increasing both limits added latency and produced slightly lower NDCG@10 than increasing only the reranker limit in each domain (Table~\ref{tab:long} and full grid in the archive).
